\documentclass[FGBook.tex]{subfiles}

\begin{document}\changefontsize{8pt}

\proza{Krysař}{VIKTOR DYK}

\noindent\begin{wrapfigure}{r}{0.4\textwidth}
\tiny

\subwection{Ukázka z díla:}
\setlength{\parindent}{3pt}
\textit{Krysař vede lid z Hammeln do propasti na hoře Koppel:}\\
Lid z Hammeln šel za~svým chmurným a~zamyšleným vůdcem, \texthl{za~velitelskou a~mocnou písní jeho píšťaly}.
Strhovala svým zvukem a~vábila dál.
Dál, až~k~bráně...\\
Klopýtali cestou necestou, padali, aby vstali a~šli dále.
Slyšeli jen krysaře a~viděli jen krysaře.
A~jak kráčeli, náhle se~před nimi otevřela propast.
Krysař stál nad propastí, chmurnější, velebnější.
Mocně a~slavně zněla píšťala.\\
Nikdo z~Hammeln se před branou nezastavil.
Nikdo, kromě Seppa Jörgena, zamyšleného nad hladinou řeky.
Všichni prošli za zvuků píšťaly s~důvěrným úsměvem a~rozepjatou náručí:
Oh, sedmihradská země!
Všichni prošli a mizeli bez hlesu.
A~chmurný a~němý vůdce pískal a~pískal.
Dav řídl za zvuků krysařovy písně; posléze krysař nad propastí osaměl -\\
Vzpomněl na \texthl{\enquote{ano}} řečené jarního večera.
Vzpomněl si na Agnes, která ho předešla, ale kterou bylo možno dohoniti.
Naklonil se nad propastí.\\
Bylo ticho, zvláštní ticho.
Píšťala vypadla z~krysařových rukou.
\texthl{Jeho píšťala znamenala život.}\\
Její zvuky jako by ještě doznívaly v pádu: její zvuky vedly krysaře, jako vedly dav -\\
"Ano," odpovídal krysař němé propasti.
A~také on~hledal bránu.\\
Tak odešel krysař a~ostatní z~Hammeln; není však jisté, došli-li do země sedmihradské.
\end{wrapfigure}

\wection{LITERÁRNÍ TEORIE}

\subwection{Literární druh a žánr:}
\noindent epická próza, novela

\subwection{Literární směr:}
\noindent\textbfhl{Literární moderna} -- vyskytuje se zde převážně \textit{impresionismus}, \textit{dekadence},
\textit{symbolismus} a~\textit{realismus}.

\subwection{Jazyk a slovní zásoba:}
\noindent Přelom \textit{19. a 20. století}, buřiči psali spisovným jazykem i~když ne tak květnatým jako modernisté.
Personifikace, metafory, řečnické otázky, citově zabarvené výrazy, přímá řeč.

% \subwection{Figury:}

% \subwection{Tropy:}

\subwection{Stylistická charakteristika textu:}
\noindent Dílo je psáno spisovným jazykem, i~když prostším.
Poetické vyznění, využití básnických figur.
Autor zde využívá vyprávěcího, popisného a~úvahového stylu.

\subwection{O díle}
\noindent Původním názvem \textbf{Pravdivý příběh}, pod kterým vycházel v časopise \textithl{Lumír}.
Knižní verze je o dvanáct let mladší a přepracovanou verzí.
Krysař vyšel i jako muzikál od Daniela Landy. \\
Námětem je jistě legenda o Krysaři, ze 13. století.
I když tato legenda pojednává o trochu jiné zápletce.
Tady Krysař z Hammeln odvede všechny děti.

\subwection{Postavy:}
\noindent 
\textbfhl{KRYSAŘ --} samozvaný nikdo; tulák; muž \textithl{bez minulosti, bez budoucnosti}.
Přichází do Hammeln jako pocestný, bere zakázku na \texthl{zbavení města krys}.
Měšťané jsou k němu zdrženliví, je nový a trochu zvláštní.
Jedinou výjimkou je mladičká Ágnes, která se do pocestného zamilovává, stejně jako on do ní.
Je to Krysař, který vede všechny občany Hammeln jako prosté krysy \textbf{vstříc jejich zkáze}.\\
\textbfhl{ÁGNES --} hlavní část \texthl{milostného trojúhelníku}, na který se \textit{Dyk} ve svém díle soustředil.
Mladá, vášnivá, krásná, nadějná.
Dobrá partie pro kohokoli ve městě, to je i důvod, proč je zasnoubená za Kristiána, dědice bohatého kupce.
A je to i důvod, proč se do ní zamiluje i takový vlk samotář, jakým je Krysař.
A ona se zamiluje do Krysaře, i když \textit{není schopna zrušit zásnuby}.
Jako jediná je otevřené mysli, nebojí se Krysaře, nebojí se venkovního světa.\\
\textbfhl{KRISTIÁN --} snoubenec Ágnes, zvaný dlouhý Kristián, pohledný mladý muž, navíc dědic velkého mění -- sám představuje dobrou partii.
S Krysařem se pouze míjí, když má ale pocit, že Ágnes ztrácí, pojistí si jí \textbf{dítětem}.\\
\textbfhl{RADNÍ --} dvojice mužů, \textithl{Gottlieb Frosch} a \textithl{Bonifác Strumm}.
Jsou zosobněním měšťanů obchodního města Hammeln.
Jestli občan nedůvěřuje cizinci, oni jím prakticky opovrhují.
Jestli občan myslí jen na peníze, jim se o nich zdají \textbfhl{sny} celé noci.
Oni, jménem celého města, najmou Krysaře na zbavení města krys.
A jsou to oni, kteří mu, jménem celého města, \textbf{odmítnou vyplatit} slíbenou odměnu.
\textbfhl{Oni představují krysy}, zkažené měšťáky a oni představují své město.\\
\textbfhl{MAGISTR --} \textit{faustovský motiv}, perfektně věří svým sebeklamům a věří, že za jeho moc může jeho mocný pán.
Nabádá Krysaře k tomu, aby využil plné síly své píšťaly na krysy.\\
\textbfhl{SEPP J\"{O}RGEN --} asi nejzáhadnější postava knihy.
Pomalu chápe, vše mu dojde až následující den, moc mu to nemyslí, lidu Hammeln je pro smích.
Děvčata si z něho utahují, lidé ho ignorují či urážejí.
A je také jediným, kdo unikne Krysařově hněvu, kdy namísto následování jeho vábení se rozhodne \texthl{postarat o plačící nemluvně}.

\subwection{Děj:}
\noindent Krysař vstupuje do města, aby ho zbavil krys.
Zamilovává se do krásné Ágnes.
Zbaví město krys a nedostává zaplaceno, stejně ale plánuje odejít.
Vrací se kvůli Ágnes, \texthl{nedokáže jí opustit}.
Potkává Magistra Faustuse, který ho nabádá k využití píšťali.
Ágnes si užívá románku s Krysařem, dokud si jí Kristián nevezme za svou.
Než mu porodit dítě, vezme si mladičká Ágnes život.
Krysaře tato zpráva zdrtí a rozhodne se zbavit \texthl{město krys}, tentokrát opravdu \texthl{všech}.
Zapíská mocně na svojí píšťalu a krysy jdou, odcházejí ze svých příbytků, zanechávají rozdělané práce, berou svoje děti a jdou.
Všechny Hammelské krysy táhnou za písní Krysařovi píšťaly.

\subwection{Kompozice, prostor a čas:}
\noindent 
Příběh se odehrává v hanzovním městě Hammeln, legenda kterou se \textbfhl{Dyk} inspiroval je datována do 13. století, ale století ve knize určené není.
Dílo sestává z 26 kapitol, chronologicky poskládaných, s občasnými zmínkami o událostech minulých.

\subwection{Význam sdělení:}
\noindent 
Linka o Krysaři se dá chápat snáze.
Město s prohnilími lidmi dostane, co si zaslouží.
Kritika povrchního měšťáctví.\\
Zajímavější je linka, která tento základní příběh upozadila.
Milostný trojúhelník Krysaře, Ágnes a dlouhého Kristiána.
Příběh muže, který nevěřil, že~může milovat a~že~může být milován.
Příběh ženy, která si~nemohla vybrat.
A~příběh muže, co~si~myslel, že~si~vybrat může.
Krysař si stále nalhával, že mu na Ágnes nezáleží, že je mu jedno, jak a~s~kým skončí, že může odejít.
A~v~tomto byl stejný jako Magistr Faustus, ne v~pekelné moci, ne v~alianci s~peklem, ale v~sebeklamu.

\wection{LITERÁRNÍ HISTORIE}

% \subwection{Politická situace:}

% \subwection{Základní principy fungování společnosti:}

% \subwection{Kontext dalších druhů umění:}

% \subwection{Kontext literárního vývoje:}

\subwection{Viktor Dyk {\ssmall -- život autora:}}
\noindent Viktor Dyk byl významným prozaikem a básníkem první poloviny 20. století.
V mládí součástí \textit{anarchistických \texthl{buřičů}}, přičinil se i ke časopisu \textit{Lumír}.
Později v životě se naklonil spíše k nacionalistickému pohledu na svět a stal se z něj i významný politik.

% \subwection{Vlivy na dané dílo:}

\subwection{Další autorova tvorba:}
\noindent 
Mezi nejznámnější díla patří \textsc{Zmoudření Dona Quijota}, \textsc{Devátá vlna} a \textsc{Od brány pekelné}.

% \subwection{Jak dílo inspirovalo další vývoj literatury:}

% \wection{LITERÁRNÍ KRITIKA}

% \subwection{Dobové vnímání díla a jeho proměny:}

% \subwection{Aktuálnost tématu a zpracování tématu:}


\wection{OSOBNÍ NÁZOR}
\noindent 
\textit{Viktor Dyk} vzal starou legendu, přikrášlil ji pro moderní literaturu a přidal k ní milostný trojúhelník.
Je to metoda, která funguje, povídka je~čtivá a~aby~nebyl konec tak předvídatelný, tak Krysařův hrůzný čin povýšil ještě na další úroveň.
V celém mém počinu jsem se nezmínil o vnitřních dialozích postav a trefně položených řečnických otázkách vypravěče/spisovatele a přitom toto jsou za mě přednosti knihy.
Ano, Dykovu novelu mohu jen doporučit.

\vfill

\noindent\begin{minipage}{\textwidth}
    {\textcolor{\wpagecolor}{\rule{\linewidth}{0.4pt}}
    \footnotesize
    \textbfhl{Legenda --} epický literární žánr pojednávající o životě mučedníka;
    \textbfhl{Literární moderna --} literární směr konce 19. století -- vyznačuje se hlavně odstupem od realismu a důrazu na jedince;
    \textbfhl{impresionismus --} požitek z daného okamžiku;
    \textbfhl{dekadence --} úpadek, rozpad;
    \textbfhl{symbolismus --} abstraktní myšlení přes skryté významy;
    \textbfhl{realismus --} věcnost bytí, správné poznání věcí samých;
    \textbfhl{Lumír --} časopis snažící se povznést českou literaturu o překlady zahraničních děl i vydávání děl českých zahraničím inspirovaných;
    \textbfhl{buřiči --} \enquote{básníci života a vzdoru}, literáti inspirovaní filosofií \texthl{Friedricha Nietzscheho}, opovrhovali společností, jejím konzervatismem a uzavřeností.
    }
\end{minipage}

\newpage

\end{document}